\ifdefined\MyWebBook\else
\documentclass[../textbook.tex]{subfiles}
\begin{document}
\fi
\bookchapter{混合专家模型}{ch:rev-moe}

稠密前馈层对每个词元应用同一组参数。增大前馈宽度可以增加表达能力，但也使每个词元承担更多矩阵计算。混合专家模型试图部分分离参数容量与单次计算：保存多组可学习变换，由输入决定本次调用其中哪些变换。减少参与计算的参数与减少全部权重的存储是两件事；跨设备通信与执行尾部另受路由结果支配。

路由器选择专家并确定混合权重，分发与合并操作将词元送入相应专家，再把结果汇回原位置。容量约束规定每个专家可接收的输入数量，溢出策略会影响输出与梯度。共享专家、软专家混合和专家并行则分别调整公共计算、输入组织和设备放置方式。

\section{条件计算与专家结构}

\subsection{专家函数的定义}
\term{混合专家模型}{Mixture of Experts}{MoE}包含若干专家函数 $F_e$，并由\term{路由器}{Router}{}或\term{门控网络}{Gating Network}{}确定各专家对当前输入的贡献。专家通常是结构相同、参数独立的前馈网络，但“结构相同”说的是结构形式，各专家训练后的函数可以任意不同。名称中的专家描述条件化计算组件，编号则只是参数索引，与数学、法律或编程这类学科分工没有对应关系。
\footnote{专家编号是模型内部的参数索引。对专家参数及路由输出通道作一致置换，可以表示同一个函数，因此编号本身不具有稳定的外部语义。专家是否形成可解释分工，需要独立的行为与干预证据。}

在 Transformer 中，常见做法是用 MoE 替换部分逐位置前馈子层，注意力和残差主干仍保留。稀疏门控使一次输入只调用少量专家，是\term{条件计算}{Conditional Computation}{}的一种形式\cite{shazeer2017sparsemoe}。它与为每种任务部署一个完整语言模型、或多个独立模型生成答案后再投票，都是不同的系统形态。

设一个专家采用门控前馈形式，行向量输入 $\vect{x}\in\Real^{1\times d}$ 的输出为
\begin{equation}
 F_e(\vect{x})=\left[\phi(\vect{x}\mat{W}_{e,g})\odot(\vect{x}\mat{W}_{e,u})\right]\mat{W}_{e,d},
\end{equation}
其中 $\mat{W}_{e,g},\mat{W}_{e,u}\in\Real^{d\times d_f}$，$\mat{W}_{e,d}\in\Real^{d_f\times d}$，$\phi$ 是固定的激活函数。每个专家都输出宽度 $d$，以便加权合并并与残差相加。普通两层前馈专家也可使用相同路由机制；专家内部结构与路由结构是两个可分别改变的维度。

\subsection{统一符号与适用假设}
将当前路由组内的有效词元表示展平为 $\mat{X}\in\Real^{N\times d}$。这里 $N$ 是实际参与路由的位置数，可以来自多个序列，但不把填充位置算入其中。路由组可能只覆盖一个本地微批次，也可能覆盖一个专家并行组；统计量必须明确采用哪一种范围。
\begin{table}[htbp]
\centering
\caption{混合专家模型核心符号与记号约定}
\label{tab:rev-moe-symbols}
\begin{tabularx}{\linewidth}{lXl}
\toprule 符号 & 含义 & 形状或范围\\\midrule
$N,d,d_f$ & 有效词元数、模型宽度、专家中间宽度 & 正整数\\
$E_r,k$ & 路由专家数、每词元所选专家数 & $1\le k\le E_r$\\
$\mat{W}_r,b_r$ & 路由器权重和偏置 & $d\times E_r,\ E_r$\\
$\mat{G},\mat{P}$ & 路由分数与全专家概率 & $N\times E_r$\\
$S_t$ & 词元 $t$ 选中的专家集合 & 大小为 $k$\\
$a_{te}$ & 稀疏混合权重 & 非负标量\\
$n_e,C_e$ & 专家申请负载与容量 & 非负整数\\
$\mat{D}_e$ & 词元到专家槽位的选择矩阵 & $C_e\times N$\\
$\mat{H}_e,\mat{Y}$ & 专家输出与合并结果 & $C_e\times d,\ N\times d$\\
$f_e,\bar p_e$ & 硬分配比例与平均路由概率 & 无量纲\\
$P_{\mathrm{total}},P_{\mathrm{active}}$ & 总参数量、每词元激活参数量 & 参数个数\\
\bottomrule
\end{tabularx}
\end{table}
本章用 $E_r$ 表示专家数量，避免与嵌入矩阵混用。$P$ 在路由部分表示专家概率矩阵，带下标的 $P_{\mathrm{total}}$ 等表示参数计数。

\begin{assumption}[基本稀疏路由模型]
路由组含 $N>0$ 个有效词元；各专家输入输出宽度相同；每个词元先选择 $k$ 个不同专家；分数有限且并列选择规则明确。先讨论没有容量溢出的情况，再显式加入容量与拒绝规则。对梯度进行局部推导时，假定选择集合和接收集合保持不变。
\end{assumption}

\subsection{稠密混合与稀疏混合}
路由器先计算
\begin{equation}
 \mat{G}=\mat{X}\mat{W}_r+\mathbf1_Nb_r^{\mathrm{T}},\qquad
 p_{te}=\frac{\exp(g_{te})}{\sum_{j=1}^{E_r}\exp(g_{tj})}.
 \label{eq:moe-router}
\end{equation}
Softmax 沿专家轴归一化。\term{稠密专家混合}{Dense Mixture of Experts}{}对每个词元执行所有专家，再计算
\begin{equation}
 \vect{y}_t=\sum_{e=1}^{E_r}p_{te}F_e(\vect{x}_t).
\end{equation}
概率可以接近零，但除非执行过程真正跳过相应专家，其算术开销仍然存在。“权重很小”不构成“没有执行”的判据。

\term{稀疏专家混合}{Sparse Mixture of Experts}{}只执行选中集合 $S_t$ 中的专家：
\begin{equation}
 S_t=\operatorname{TopK}(g_{t,:},k),\qquad
 \vect{y}_t=\sum_{e\in S_t}a_{te}F_e(\vect{x}_t).
 \label{eq:moe-sparse}
\end{equation}
Top-$k$ 表示取最高的 $k$ 个分数对应的不同编号。由于 Softmax 对同一行保持分数次序，在没有额外变换时，按分数或概率选择得到同一集合。概率如何用于合并，则另有不同约定。

\section{路由权重与局部梯度}

\subsection{两种归一化约定}
一种常见形式在选中专家内重新归一化：
\begin{equation}
 a_{te}=\begin{cases}
 \displaystyle\frac{\exp(g_{te})}{\sum_{j\in S_t}\exp(g_{tj})},&e\in S_t,\\
 0,&e\notin S_t.
 \end{cases}
 \label{eq:moe-selected-normalization}
\end{equation}
此时每个词元保留的混合权重和为一。也可以先对全部专家归一化，再直接保留选中的概率，即 $a_{te}=p_{te}\mathbf1[e\in S_t]$。后者保留权重之和一般小于一；输出幅度因而也受到路由置信度影响。Switch 的单专家路由保留了被选专家的全专家门控概率\cite{fedus2022switch}，选中集合内重新归一化会得到另一个函数。

两种形式在 $k=1$ 时差异尤其明显。选中集合内归一化会得到唯一权重为一，在固定选择集合的邻域内，任务损失无法再通过这个恒为一的权重更新路由分数。保留全专家概率则仍有连续梯度路径。判断路由器能否从任务损失学习，依据是后续归一化是否消去了 Softmax 的作用，代码中出现 Softmax 本身并不说明问题。

并列分数处需要一个明确的编号选择规则；选择边界附近，微小浮点扰动也可能改变专家集合。确定性并列规则可以使重复执行更可解释，但不消除硬选择本身的非光滑性。

\subsection{专家与路由器梯度}
暂时固定一个词元的选中集合，并省略词元下标。令专家输出 $\vect{z}_e=F_e(\vect{x})$，混合结果 $\vect{y}=\sum_ea_e\vect{z}_e$，上游梯度为列向量 $\vect{h}=\frac{\partial\mathcal L}{\partial \vect{y}}$。定义标量敏感度 $r_e=\vect{h}^{\mathrm{T}}\vect{z}_e^{\mathrm{T}}$，则
\begin{equation}
 \frac{\partial\mathcal L}{\partial a_e}=r_e,\qquad
 \frac{\partial\mathcal L}{\partial \vect{z}_e}=a_e\vect{h}^{\mathrm{T}}.
\end{equation}
专家参数的梯度由其局部网络反向传播获得，并乘以该词元的混合权重。未选专家不接收这个词元经专家执行路径产生的梯度，但可能从同批其他词元、正则项或其他共享路径获得贡献。

对式\eqref{eq:moe-selected-normalization}，选中集合内的 Softmax 雅可比为 $a_e(\mathbf1[e=j]-a_j)$，所以
\begin{align}
 \frac{\partial\mathcal L}{\partial g_j}
 &=\sum_{e\in S}r_ea_e(\mathbf1[e=j]-a_j)\\
 &=a_j\left(r_j-\sum_{e\in S}a_er_e\right),\qquad j\in S.
 \label{eq:moe-router-grad}
\end{align}
未选分数的这一局部梯度为零。式中括号比较某个专家的输出敏感度与当前混合平均值；因此路由器可以学习改变已选专家之间的连续权重。

若保留全专家概率，令 $r_e=0$ 对所有未选专家成立，则同样的 Softmax 反向得到
\begin{equation}
 \frac{\partial\mathcal L}{\partial g_j}
 =p_j\left(r_j-\sum_{e\in S}p_er_e\right),\qquad j=1,\ldots,E_r.
\end{equation}
全局归一化使未选专家的路由分数也可获得梯度，而专家参数的任务梯度仍来自实际执行的路径。

最后，若所有词元的分数梯度构成 $\mat{G}_G$，路由器线性层满足
\begin{equation}
 \frac{\partial\mathcal L}{\partial \mat{W}_r}=\mat{X}^{\mathrm{T}}\mat{G}_G,\qquad
 \frac{\partial\mathcal L}{\partial b_r}=\mat{G}_G^{\mathrm{T}}\mathbf1_N.
\end{equation}
输入梯度还应加上所有已选专家的回传与路由器的回传。漏掉任一路径，得到的就不再是完整混合函数的导数。

\subsection{硬选择的导数边界}
Top-$k$ 的编号集合在分数排序不变的小邻域中保持常数，在排序交换处发生跳变。通常的反向传播沿已选择路径求导，不会自动告诉模型“如果换成另一个专家，损失会怎样”。加噪声、辅助目标或其他选择机制可以改变探索和学习行为，但每种方法都需说明自己的前向与反向定义。

如果所有词元过早集中到少数专家，这些专家得到更多任务训练机会，其他专家则缺少有效更新，可能形成自我强化的负载差异。负载均衡机制因此不仅是运行调度，也影响训练过程中不同参数获得监督的分布。均衡目标还需与专家分工和任务质量相协调。

\section{词元路由}

\subsection{词元轴的专家批次变换}
MoE 的输出仍按原词元顺序排列，但专家计算希望把送往同一专家的词元放在同一个批次中。\term{分发}{Dispatch}{}将原词元重排或复制到专家局部缓冲，\term{合并}{Combine}{}则将处理后的结果按原位置加权累加。

专家 $e$ 有 $C_e$ 个槽位，定义选择矩阵 $\mat{D}_e\in\{0,1\}^{C_e\times N}$。如果原词元 $t$ 被接收到槽位 $c$，令 $(\mat{D}_e)_{ct}=1$；一个槽位至多接收一个词元。于是
\begin{equation}
 \mat{X}_e=\mat{D}_e\mat{X},\qquad \mat{H}_e=F_e(\mat{X}_e),
\end{equation}
其中 $F_e$ 对每个有效槽位按行计算。再定义合并矩阵 $\mat{A}_e\in\Real^{N\times C_e}$，将对应位置设为接收后的权重，得到
\begin{equation}
 \mat{Y}=\sum_{e=1}^{E_r}\mat{A}_e\mat{H}_e.
 \label{eq:moe-dispatch-combine}
\end{equation}
没有拒绝且采用原权重时，$(\mat{A}_e)_{tc}=a_{te}(\mat{D}_e)_{ct}$。这一写法用于说明映射关系，不要求运行时物化巨大稀疏矩阵；索引列表、前缀和与分组缓冲可以表达同样的关系。

图\ref{fig:moe-dispatch}概括了分组与回收路径。一个词元选择两个专家时，分发阶段逻辑上产生两份输入引用，合并时必须将两份结果相加。只保留最后返回结果会丢失一条专家路径。空槽位即使因偏置产生非零专家输出，其合并权重也必须为零；否则填充缓冲会污染真实词元。

\begin{figure}[htbp]
\centering
\begin{tikzpicture}[>=Stealth,every node/.style={font=\small},
 box/.style={draw,rounded corners,align=center,minimum width=22mm,minimum height=10mm}]
\node[box] (x) at (0,0) {词元表示\\$N\times d$};
\node[box] (r) at (2.8,0) {路由与分发\\编号、权重、槽位};
\node[box] (e1) at (6,1.3) {专家1\\$n_1\times d$};
\node[box] (e2) at (6,0) {专家2\\$n_2\times d$};
\node[box] (e3) at (6,-1.3) {专家3\\$n_3\times d$};
\node[box] (y) at (9.4,0) {回传与加权合并\\$N\times d$};
\draw[->] (x)--(r);
\foreach \e in {e1,e2,e3}{\draw[->] (r)--(\e);\draw[->] (\e)--(y);}
\node[align=center] at (4.8,-2.35) {词元顺序转为专家分组顺序，再恢复原位置；\\专家可以位于不同设备，箭头可能跨越网络。};
\end{tikzpicture}
\caption{稀疏 MoE 的数据流。专家批次大小由路由决定；选择多个专家意味着多条返回路径。}
\label{fig:moe-dispatch}
\end{figure}

\subsection{路由容量与溢出}
令申请负载 $n_e=\sum_t\mathbf1[e\in S_t]$，无拒绝时有 $\sum_en_e=Nk$，平均负载为 $\frac{Nk}{E_r}$。一种容量规则是
\begin{equation}
 C=\left\lceil\gamma\frac{Nk}{E_r}\right\rceil,
 \label{eq:moe-capacity}
\end{equation}
其中 $\gamma>0$ 为\term{容量因子}{Capacity Factor}{}，所有专家暂取相同容量。Switch 中的单专家形式对应 $k=1$\cite{fedus2022switch}；其他系统可能按局部分组、不同专家宽度或不同对齐规则定义容量，容量因子的数字只有在这些约定一致时才可比。

即使总槽位 $E_rC\ge Nk$，局部专家仍可能溢出。例如一个专家收到大量申请时，其他专家的空槽位不能自动代替它执行同一个函数。容量通常与静态缓冲大小、通信消息形状和最坏负载相关，提高容量可以减少拒绝，却可能增加填充计算和临时空间。

超过容量后的处理必须成为模型定义的一部分。拒绝某条专家路径，可以使该路径贡献为零；重新路由会调用另一组参数；等待或变长分组可以保留全部路径，但付出动态调度和尾延迟代价。被称为“不丢词元”的实现通常是保留全部路由计算，它替换掉的是丢弃路径这一策略，负载和通信约束依然存在。

若一个词元的全部专家路径被拒绝，MoE 子层贡献可能为零，残差主干仍可保留该位置。这里的“丢词元”常指跳过该层的专家处理，序列中的编号和它的监督标签都不受这一操作影响。
\footnote{不同论文与运行时对 token dropping 的统计口径可能不同：拒绝的专家分配数、至少失去一条路径的词元数、以及失去全部路径的词元数，应分别报告。}

\subsection{接收规则可能引入批次依赖}
容量竞争发生在路由组内。若采用先到先得，后部位置可能更容易被拒绝；若按门控权重排序，则低权重路径更容易失去计算。于是同一词元在不同批次组成中，可能因为其他词元的申请而获得不同输出。这样的函数依赖同批其他词元，与逐词元独立的前馈函数不同。

对于严格自回归训练，还应警惕一种间接依赖：若先观察整段序列的路由权重，再按全局排名决定早期词元能否获得容量，未来词元可能影响早期位置的接收结果。避免这一问题需要容量方案和接收规则与因果定义相容，注意力层的因果掩码对此无能为力。实际设计应区分前向内容读取、容量竞争以及训练辅助统计三类关系。

\section{负载均衡与路由稳定性}

\subsection{硬负载与软概率描述不同对象}
定义容量截断前的分配比例和平均概率：
\begin{equation}
 f_e=\frac1{Nk}\sum_{t=1}^{N}\mathbf1[e\in S_t],\qquad
 \bar p_e=\frac1N\sum_{t=1}^{N}p_{te}.
 \label{eq:moe-load-stats}
\end{equation}
二者各自之和为一。$f_e$ 衡量实际申请的离散次数，$\bar p_e$ 衡量路由器的连续概率质量。一个专家经常以第二名被选中，可能有较大的申请次数，却未必拥有最大的平均概率，因此两种统计各有各的用途。

若只统计容量截断后的负载，被饱和容量裁平的结果可能看起来很均衡，却掩盖大量溢出。应同时保留申请数、接收数、拒绝数和执行耗时。若专家宽度不同，次数均衡也不等于算术量均衡。

\subsection{辅助均衡目标与梯度范围}
Switch 使用硬分配比例与平均概率乘积构造辅助目标。按照式\eqref{eq:moe-load-stats}的 top-$k$ 归一化，本章讨论其一种相应扩展：
\begin{equation}
 \mathcal L_{\mathrm{bal}}=E_r\sum_{e=1}^{E_r}f_e\bar p_e,\qquad
 \mathcal L=\mathcal L_{\mathrm{task}}+\lambda\mathcal L_{\mathrm{bal}}.
 \label{eq:moe-loadloss}
\end{equation}
其中 $\lambda\ge0$ 控制任务目标与均衡目标的相对权重，top-1 时恢复相应统计定义。不同实现对 $k$、序列和设备的归一化可能不同，系数不能脱离公式移植。

在反向时通常将硬计数 $f_e$ 视为当前已确定的常数，而通过 $\bar p_e$ 对路由器求导。由 Softmax 雅可比，
\begin{align}
 \frac{\partial\mathcal L_{\mathrm{bal}}}{\partial g_{tj}}
 &=\frac{E_r}{N}\sum_ef_ep_{te}(\mathbf1[e=j]-p_{tj})\\
 &=\frac{E_r}{N}p_{tj}\left(f_j-\sum_ef_ep_{te}\right).
 \label{eq:moe-balance-grad}
\end{align}
因此它可以为未被选择的分数提供连续训练信号，离散计数本身则没有参与求导。它是实际最大执行时间之外的替代指标，每一步更新后负载是否立刻均匀另需观察。

当 $f_e=\bar p_e=\frac{1}{E_r}$ 时，该辅助项为1，形成一个便于比较的均匀参考值。它衡量的是这一乘积而非均衡偏差的范数，某个标量接近1也不足以认定系统均衡。特别是硬分配与软概率的关系、局部分组及有限样本都会影响这个乘积；应联合观察各专家分布和最大负载。

\subsection{统计范围与任务权重}
如果各设备本地批量大小不同，先计算各自平均概率再无权平均，会改变全局统计。应先合并概率和与有效词元数，再形成所需范围内的均值。另一方面，全局均衡可能掩盖单台设备上的局部热点；设计目标必须说明是专家层级、设备层级还是通信组层级。

均衡压力过大时，路由可能优先满足平均分配而牺牲任务有效性；压力过小时，少数专家可能成为瓶颈并获得不成比例的更新。除了辅助损失，也可以通过容量分配、专家复制或路由偏置控制负载，各策略机制对比如表\ref{tab:rev-moe-balance-strategies}所示。

\begin{table}[htbp]
\centering
\caption{主流 MoE 负载均衡与路由正则化策略对比}
\label{tab:rev-moe-balance-strategies}
\begin{tabularx}{\linewidth}{lXXX}
\toprule
均衡策略 & 作用机制与公式 & 核心优势 & 局限与权衡\\\midrule
辅助损失（Aux-Loss） & $\mathcal{L}_{\mathrm{bal}} = E_r \sum_e f_e \bar{p}_e$ 加入总损失 & 梯度显式驱动路由器分布均匀 & 需调节平衡权重 $\lambda$，易损害语言建模主目标\\
无辅助损失偏置 & 动态调整专家打分偏置 $g_e \gets g_e + b_e$ & 消除辅助损失对主干梯度的干扰 & 需基于设备负载滚动自适应微调偏置步长\\
容量因子与丢弃 & 设定每专家容量上限 $C = \lceil\gamma \frac{Nk}{E_r}\rceil$ & 严格固定通信与计算缓冲区大小 & 低因子下词元丢弃损害序列语义完整性\\
专家选择（Expert Choice） & 专家主动挑选前 $k$ 个最高分词元 & 自然消除专家负载不均与溢出 & 词元被选次数不均，部分词元可能未被分配\\
共享专家解耦 & 常驻独立共享专家捕获公共表征 & 路由专家专注特定差异化特征 & 增加每词元固定计算底噪，无法稀疏跳过\\
\bottomrule
\end{tabularx}
\end{table}

路由分数的数值范围同样影响稳定性。稳定 Softmax 通过减去行最大值避免指数溢出，分数排序因低精度舍入而变化则不在它的覆盖范围内。路由计算与专家矩阵乘法可以采用不同精度；是否合适应结合选择差异、梯度和运行成本判断，全部运算统一数据类型并不是必要条件。

\begin{example}
取 $N=d=4$，$E_r=3$，每词元选择 $k=2$ 个专家，使用选中集合内重新归一化。令输入 $\mat{X}=I_4$，路由器无偏置，并设
\begin{equation}
 \mat{W}_r=\begin{bmatrix}
 \log0.7&\log0.2&\log0.1\\
 \log0.6&\log0.3&\log0.1\\
 \log0.55&\log0.35&\log0.1\\
 \log0.1&\log0.2&\log0.7
 \end{bmatrix}.
\end{equation}
因为每行概率之和为一，所以 $\mat{G}=\mat{X}\mat{W}_r=\mat{W}_r$，Softmax 后恰好得到矩阵中的对数自变量。设三个专家为线性映射 $F_1(\vect{x})=\vect{x}$、$F_2(\vect{x})=2\vect{x}$、$F_3(\vect{x})=3\vect{x}$。各词元分配细节如表\ref{tab:rev-moe-example-routing}所示。

\begin{table}[htbp]
\centering
\caption{四词元在三专家路由下的概率分布、激活专家与选中权重}
\label{tab:rev-moe-example-routing}
\begin{tabular}{clll}
\toprule 词元 & 全专家概率 & 选中集合 & 选中后权重\\\midrule
1 & $(0.7,0.2,0.1)$ & $\{1,2\}$ & $(\frac{7}{9},\frac{2}{9})$\\
2 & $(0.6,0.3,0.1)$ & $\{1,2\}$ & $(\frac{2}{3},\frac{1}{3})$\\
3 & $(0.55,0.35,0.1)$ & $\{1,2\}$ & $(\frac{11}{18},\frac{7}{18})$\\
4 & $(0.1,0.2,0.7)$ & $\{2,3\}$ & $(\frac{2}{9},\frac{7}{9})$\\
\bottomrule
\end{tabular}
\end{table}
最后一行按专家编号列权重，与按分数从高到低排列得到的是不同顺序。保存路由记录时必须同时保存编号与权重，两种排序未必一致。

没有容量拒绝时，四个输出分别为
\begin{align}
 \vect{y}_1&=(\frac{7}{9})\vect{x}_1+(\frac{2}{9})(2\vect{x}_1)=(\frac{11}{9})\vect{x}_1,\\
 \vect{y}_2&=(\frac{2}{3})\vect{x}_2+(\frac{1}{3})(2\vect{x}_2)=(\frac{4}{3})\vect{x}_2,\\
 \vect{y}_3&=(\frac{11}{18})\vect{x}_3+(\frac{7}{18})(2\vect{x}_3)=(\frac{25}{18})\vect{x}_3,\\
 \vect{y}_4&=(\frac{2}{9})(2\vect{x}_4)+(\frac{7}{9})(3\vect{x}_4)=(\frac{25}{9})\vect{x}_4.
\end{align}
由于 $\mat{X}=I_4$，$\mat{Y}$ 就是以这些系数为对角元素的矩阵。路由权重和为一，但输出系数可以大于一，因为专家输出本身并不被限制在单位范围。

专家申请数为 $(n_1,n_2,n_3)=(3,4,1)$，总计8条路径。取容量因子 $\gamma=1$，由式\eqref{eq:moe-capacity}得到每专家容量 $C=3$。虽然总槽位数为9，大于申请总数8，专家2仍会溢出。

设按词元顺序接收，专家2接收词元1、2、3，拒绝词元4。若拒绝路径直接置零且不重新归一化，词元4输出成为
\begin{equation}
 \vect{y}_4^{\mathrm{drop}}=(\frac{7}{9})(3\vect{x}_4)=(\frac{7}{3})\vect{x}_4.
\end{equation}
若改为对剩余接收路径重新归一化，唯一专家3的权重变为一，输出则为 $3\vect{x}_4$。两个结果均不同于原来的 $(\frac{25}{9})\vect{x}_4$，而且彼此也不相等。溢出策略因此造成明确的函数差异，它不只是不影响结果的内存优化。

本例拒绝1条路径，路径拒绝率为 $\frac{1}{8}$；有1个词元失去部分路径，词元受影响比例为 $\frac{1}{4}$；没有词元失去全部路径，完全跳过专家层的比例为零。这三个数字分别回答不同问题。

截断前的归一化分配比例为
\begin{equation}
 f=(\frac{3}{8},\frac{1}{2},\frac{1}{8}),\qquad
 \bar p=(0.4875,0.2625,0.25).
\end{equation}
代入式\eqref{eq:moe-loadloss}得到
\begin{equation}
 \mathcal L_{\mathrm{bal}}
 =3\left[(\frac{3}{8})(0.4875)+(\frac{1}{2})(0.2625)+(\frac{1}{8})(0.25)\right]
 =1.0359375.
\end{equation}
专家2的申请次数最多，却不是平均概率最高的专家，因为它在多个词元上作为第二选择进入集合。这解释了为什么仅检查平均门控概率不足以诊断实际执行负载。

最大申请负载相对平均负载为 $\frac{4}{\frac{8}{3}}=1.5$。若只观察截断后的接收数 $(3,3,1)$，这个热点会显得减弱，但那是拒绝计算造成的结果。质量分析和性能分析都必须保留截断前后的统计。
\end{example}

\begin{algorithm}[带显式容量规则的稀疏专家前向]\label{alg:moe-routing}
\AlgoInput{$\mat{X}\in\Real^{N\times d}$，路由器 $\mat{W}_r,b_r$，$E_r$ 个专家，选择数 $k$，容量 $C$；采用选中内归一化、按词元顺序接收、拒绝路径置零且不二次归一化。}
\AlgoOutput{$\mat{Y}\in\Real^{N\times d}$，申请数、接收数、拒绝数和路由记录。}
\AlgoState{每专家槽位计数 $c_e=0$，输出 $\mat{Y}=0$，以及接收记录 $(t,e,c,a_{te})$。}
\begin{enumerate}
\item 若 $N=0$，返回空输出和零计数；否则检查 $1\le k\le E_r$、$C\ge1$，计算路由分数及全专家概率。非有限分数作为输入失败处理，不进入排序。
\item 对词元 $t=1,\ldots,N$，按分数降序选取不同的 $k$ 个专家；并列时按专家编号升序。计算选中集合内权重，记录所有申请。
\item 对该词元每条选中路径，若 $c_e<C$，将 $\vect{x}_t$ 写入专家 $e$ 的下一个槽位，并记录原词元编号与权重；增加 $c_e$。否则记录拒绝，不产生专家输入。
\item 每个专家仅对其接收的 $c_e$ 个有效槽位计算 $F_e$。专家位于远端时，按同一槽位记录分发输入并回收输出。
\item 对每条接收记录执行 $\mat{Y}_{t,:}\leftarrow \mat{Y}_{t,:}+a_{te}(\mat{H}_e)_{c,:}$。全部接收记录合并完毕即停止。
\end{enumerate}
\textbf{不变量：}每个槽位至多对应一个输入；每条接收记录仅合并一次；每专家接收数不超过容量；申请数等于接收数加拒绝数。若无拒绝，每词元合并权重和为一；有拒绝时本算法不再要求这一行和为一。

路由线性投影计算量为 $O(NdE_r)$；直接完整排序的上界为 $O(NE_r\log E_r)$，选择算法可避免完整排序。路由记录为 $O(Nk)$，专家计算按实际接收数计。分发缓冲、网络代价和各专家矩阵大小需另行计入。
\end{algorithm}
算法\ref{alg:moe-routing}完整定义了一种容量受限函数。修改排序、二次归一化或拒绝规则时，应先修改这个函数定义，再讨论底层实现是否保持相同结果。它并不要求高性能运行时逐词元串行执行，批量排序和索引操作可以在保持接收语义的前提下并行化。

\section{总参数、激活参数与执行成本}

\subsection{参数容量与词元计算}
设非路由部分参数为 $P_0$，每个路由专家有 $P_e$ 个参数，路由器有 $P_r$ 个参数，另有始终执行的共享专家参数 $P_s$。在各路由专家同尺寸、没有拒绝的单层计数下，
\begin{equation}
 P_{\mathrm{total}}=P_0+P_r+P_s+E_rP_e,\qquad
 P_{\mathrm{active}}=P_0+P_r+P_s+kP_e.
\end{equation}
对于多层模型应逐层累加，各层未必都是同样的 MoE。每词元激活比例 $\frac{k}{E_r}$ 仅对应可路由专家部分，整个模型的计算比例还包含其余部分。

门控前馈专家忽略偏置时约有 $P_e=3dd_f$ 个参数。若取 $d=1024$、$d_f=4096$、$E_r=8$、$k=2$，则单专家参数为 $12\,582\,912$，全部路由专家为 $100\,663\,296$，每词元激活的路由专家参数为 $25\,165\,824$。

以每权重两字节存储，全部路由专家权重为 $201\,326\,592$ 字节，即192 MiB；每词元只执行其中两组专家，常驻权重预算仍需按192 MiB计算。实际内存还包括路由器、共享层、缓存、激活和工作空间。

\begin{warning}[专家参数与驻留容量]
稀疏性减少每词元经过的专家路径；总权重与优化器状态仍由完整参数集合及其放置方式决定。专家分片或卸载可以降低单设备常驻量，但需计入通信、加载与峰值缓冲，权重预算应以完整参数集合为准。
\end{warning}

\subsection{批次专家覆盖率}
单个词元选择少量专家，不表示整个批次只读取少量权重。若多个词元选择不同专家，批次专家集合是 $\bigcup_tS_t$，其大小可以接近 $E_r$。增大批量可能改善每个专家的矩阵尺寸，也可能增加本步触及的权重范围，两种效应需要同时分析。

在乘法与加法分别计一次浮点操作、专家主要由矩阵乘法组成的近似下，执行 $N$ 个词元的专家前向成本约为 $2NkP_e$。共享部分、路由排序、激活函数、通信与填充不包含在这个式子中；训练的反向成本还取决于保存和重计算的中间状态。稠密模型的 $6P\mat{D}$ 估算代入总 MoE 参数会高估实际执行量，只代入激活参数又会漏掉共享部分与通信，两者都得不到精确时间。

\subsection{负载不均与小矩阵效率}
若每个专家处理一个局部矩阵乘法，理想均衡时每专家批量约为 $\frac{Nk}{E_r}$。专家数增加而总词元数不变，会减小这一批量。矩阵变小可能降低设备利用率，路由和内核启动开销占比则升高。增加专家数换取容量，代价正落在这里。

各设备必须等待相关专家结果返回，执行尾部受最慢专家或最慢设备影响。即使平均专家耗时很低，单个热点也可能决定整个层的完成时间。应关注最大负载、最大接收延迟和跨设备流量，而不只是所有专家平均的利用率。

\section{共享专家、细粒度专家与软混合}
\optional
\subsection{共享路径与细粒度组合}
\term{共享专家}{Shared Expert}{}始终处理每个词元，路由专家则条件执行。一种输出形式为
\begin{equation}
 \vect{y}_t=F_{\mathrm{shared}}(\vect{x}_t)+\sum_{e\in S_t}a_{te}F_e(\vect{x}_t).
\end{equation}
共享项没有被路由权重缩小，其参数与算术开销都应每次计入。共享专家提供始终执行的变换，路由专家提供随输入选择的变换；两类路径的实际分工由训练形成。

\term{细粒度专家}{Fine-Grained Expert}{}将较宽专家改为更多较窄专家，同时增加被选专家数。例如将中间宽度除以 $m$、专家数乘以 $m$、选择数也乘以 $m$，在忽略路由和其他开销时，专家总参数与每词元专家矩阵计算可以近似保持不变。这增加了组合自由度，却使路由分数、分组和通信的组织更加复杂。DeepSeekMoE 将细粒度划分与共享专家作为两项结构设计\cite{dai2024deepseekmoe}；具体型号仍应按精确配置和权重核对。

专家划分不是将一个已训练稠密矩阵任意切块后就自动得到等价 MoE。激活函数、门控组合以及条件选择会改变函数，需要相应训练或经过论证的转换过程。

\subsection{软专家混合改变专家输入}
\term{软专家混合}{Soft Mixture of Experts}{Soft MoE}让专家接收词元的连续混合，与给所有专家保留一个很小概率的稠密混合不同。其一种构造先把词元连续混合成有限个专家槽位，再由专家处理槽位，最后将槽位输出混合回词元\cite{puigcerver2024softmoe}。

设总槽位数 $S=E_rp$，每专家有 $p$ 个槽位，槽位打分矩阵为 $\mat{Z}=\mat{X}\mat{\Phi}\in\Real^{N\times S}$。分发权重对每个槽位沿词元轴归一化，合并权重对每个词元沿槽位轴归一化：
\begin{equation}
 D_{ts}=\frac{\conste{}^{Z_{ts}}}{\sum_{u=1}^{N}\conste{}^{Z_{us}}},\qquad
 C_{ts}=\frac{\conste{}^{Z_{ts}}}{\sum_{v=1}^{S}\conste{}^{Z_{tv}}}.
\end{equation}
于是槽位输入 $\widetilde{\mat{X}}=\mat{D}^{\mathrm{T}}\mat{X}\in\Real^{S\times d}$，各槽位按归属专家计算 $\widetilde{\mat{Y}}_s=F_{e(s)}(\widetilde{\mat{X}}_s)$，最后 $\mat{Y}=\mat{C}\widetilde{\mat{Y}}$。这里的 $\mat{D}$ 是连续混合权重，区别于前文硬选择矩阵 $\mat{D}_e$。

每个专家处理的是若干词元的加权组合，与对每个原始词元独立执行一次专家的稠密混合不同。两次归一化也不能交换：分发列和为一，合并行和为一，各自保证不同的混合对象。

这类全局词元混合若直接用于自回归解码，会把未来位置混入早期位置的专家输入，替换因果语言模型中的逐位置 MoE 需要另行定义因果槽位构造并分析代价。

\subsection{路由方式的约束差异}
\begin{table}[htbp]
\centering
\caption{主流专家路由方式、处理粒度与约束特性对比}
\label{tab:rev-moe-routing-modes}
\begin{tabularx}{\linewidth}{lXXX}
\toprule 方式 & 专家处理对象 & 离散选择 & 主要约束\\\midrule
稠密混合 & 每个词元 & 无需截断 & 每词元执行全部专家\\
词元选择的稀疏混合 & 选中词元 & 每词元 top-$k$ & 专家负载不固定\\
专家选择的稀疏混合 & 专家挑选的词元 & 每专家挑选输入 & 每词元路径数可变\\
软专家混合 & 连续混合的槽位 & 不必使用硬 top-$k$ & 槽位数与混合可见性\\
\bottomrule
\end{tabularx}
\end{table}
专家选择表示按专家分数列挑选有限词元，能够直接约束每专家的接收量，但一个词元可能得到多个专家，也可能没有专家处理。若需要固定每词元覆盖，还要增加额外约束。比较方法时，应分别说明固定的是词元路径数、专家容量还是槽位数量，“稀疏”一词覆盖不了这些不同定义。

\section{专家并行的通信与训练推理差异}

\subsection{参数放置与词元移动}
\term{专家并行}{Expert Parallelism}{EP}把不同专家参数放在不同设备。路由完成后，词元必须送到相应设备；专家输出再送回原来的词元归属设备。\term{全互换通信}{All-to-All}{}可以表达这种多发送端到多接收端的数据交换，各设备之间的流量由词元路由结果决定。

设远端执行的接收路径比例为 $\rho$，每激活元素占 $s_a$ 字节，输入输出宽度均为 $d$。忽略元数据和协议开销，前向分发与返回的通信量级为
\begin{equation}
 V_{\mathrm{comm}}\approx2\rho Nkd s_a.
\end{equation}
该式按每条远端路径的单向有效载荷计数，并将发送输入与返回输出各计一次；不同统计工具若同时累计发送和接收端会得到不同口径。张量并行、专家复制、压缩传输与容量拒绝都会改变实际数值。

理想传输时间至少受有效带宽与启动延迟限制。若共有 $m$ 轮不可忽略的启动，延迟为 $\ell$、有效带宽为 $\vect{b}_{\mathrm{eff}}$，可以用
\begin{equation}
 t_{\mathrm{comm}}\gtrsim m\ell+\frac{V_{\mathrm{comm}}}{\vect{b}_{\mathrm{eff}}}
\end{equation}
理解瓶颈；拥塞、同步和拓扑争用还会增加时间。计算与通信只有在依赖允许时才能重叠，理论算力与网络带宽无法同时按满负荷利用。

\subsection{专家并行与数据并行}
数据并行复制同一模型参数并处理不同数据，专家并行则让一次前向调用不同设备上的不同专家参数。专家参数只在拥有同一专家副本的数据并行组内同步梯度，路由器与共享参数又可能采用不同复制范围。组定义错误会造成不同参数被错误求和，或应保持一致的副本没有同步。

同一设备可容纳多个专家，单个专家也可进一步进行张量切分。专家编号、物理设备编号和张量分片编号不是一个索引。模型配置记录逻辑专家结构，运行时放置表记录执行位置，检查点必须能从分片还原这些关系。

\subsection{训练与推理张量形状}
训练往往同时处理较多词元，专家批次相对较大，但需保存反向状态并执行反向通信。各专家梯度还可能十分稀疏或不均匀，优化器状态则不能仅按当前激活量估计。激活重计算需要重现路由选择、随机噪声与容量接收结果，否则反向不再对应原前向路径。

预填充可以并行处理一段输入，而逐词元解码每个请求每步通常只新增一个位置。并发较小时，各专家收到的矩阵可能很小，路由、权重读取和通信启动的比例增大。增加请求合批可能提高专家利用率，却引入等待时间及更复杂的容量竞争。

训练中接受的路径拒绝比例不应直接移植到服务端。一个被拒绝的关键词元可能改变后续生成，且输出依赖批次组成会使复现和回归更困难。推理可以采用变长无拒绝执行、扩大容量或其他明确方案，每项变化都需要评估其函数、质量和延迟影响。

\section{可执行结构还原}

模型名称中出现 MoE 或激活参数数字，只能提供结构线索。需要按每层核对路由专家数、共享专家数、每词元选择数、中间宽度、门控归一化、路由激活、容量策略及哪些层使用 MoE。再将这些配置与路由矩阵、专家参数张量和前向定义相互对应。主流开源 MoE 模型核心架构与路由配置参数如表\ref{tab:rev-moe-architectures}所示。

\begin{table}[htbp]
\centering
\caption{主流开源 MoE 大语言模型架构参数与路由策略全景对比}
\label{tab:rev-moe-architectures}
\begin{tabularx}{\linewidth}{lXXXXX}
\toprule
模型 & 总参数 / 激活参数 & 专家总数 $E_r$ & 激活专家数 $k$ & 共享专家数 & 负载均衡与路由特性\\\midrule
Switch-Base & $7.4\text{B} / 0.9\text{B}$ & $128$ & $1$ & $0$ & 容量因子截断 + 辅助损失，极度稀疏\\
Mixtral 8x7B & $46.7\text{B} / 12.9\text{B}$ & $8$ & $2$ & $0$ & 粗粒度标准专家，Softmax 选中重新归一化\\
Qwen1.5-MoE & $14.3\text{B} / 2.7\text{B}$ & $64$ & $4$ & $4$ & 细粒度专家分割 + 固定共享专家协同\\
DeepSeek-V2 & $236\text{B} / 21\text{B}$ & $160$ & $6$ & $2$ & 细粒度（$0.4\times$ 宽度）+ 动态设备亲和路由\\
DeepSeek-V3 & $671\text{B} / 37\text{B}$ & $256$ & $8$ & $1$ & 无辅助损失路由（Bias-Only 平衡）+ 节点限流通信\\
\bottomrule
\end{tabularx}
\end{table}

例如路由投影在本章约定中为 $d\times E_r$；某些库按输出乘输入布局保存为 $E_r\times d$，这属于转置约定，形状方向不同并不说明其中一种有误。量化参数可能按组打包，存储元素数与原始参数计数也不相同。对于共享专家，既要确认参数存在，也要确认前向确实每次调用，张量名称不足以证明这一点。

模型分析可从路由分布、参数与配置对应关系、任务质量和运行性能四方面展开。路由分布反映输入分配，专家专业化需要通过跨样本观察与消融研究；配置和前向定义决定合并语义，设备拓扑与实际负载共同决定吞吐。






\chapterreferences
\myendchapterdocument
